\documentclass[]{../../../../tex/classes/homework}
\usepackage{../../../../tex/packages/hebrewSupport}

\usepackage{tabularx}

\author{שחר פרץ}
\title{עוד עבודה ללמידה עצמאית בהיסטוריה}
\begin{document}
	\maketitle
	\section{}
	\textbf{שאלה: }\\
	מהי תנועת המרי? כיצד היא פעלה? הצידו שתי פעולות שביצעה התנועה והסבירו כיצד פעולות אלא מבטאות את אופי הפעילות שהסברתם. 
	\textbf{תשובה: }\\
	בסוף המאה ה־19 התחלה להתגבש תנועת ה\textit{ציונות}, תנועה השואפת להקים בית מדיני לעם היהודי בארץ ישראל. התנועה החלה את פעולתה המאורגנת ע''י בינימין זאב הרצל, שהקים את רוב מסודותיה המאורגנים. בשנת 1945 נגמרה מלחמה בשם מלחמת העולם השנייה, שבה נכלאו ונעקרו מבתיהם יהודים רבים. 
	
	בא''י באותה התקופה שלטה בריטניה, ונהגה במדיניות שנקראה ''ספר הלבן השלישי``. המדיניות הגבילה באופן משמעותי את כמות היהודים שיכלו לעלות לא''י, ואת התפשטות היישוב היהודי. פעלו בארץ שלוש מחתרות ציוניות עיקריות – ההגנה, שהייתה חלק מהנציגות ה''רשמית`` של התנועה הציונית, והובילה קו פעולה שהתאפיין בתמיכה בתחום הדיפלומטיה והעלייה, ואכן האצ''ל והלח''י שהתנגדו למדיניות ''ההגנה`` ופעלו נגד הבריטים, אזרחיהם ומוסדותיהם באופן צבאי. כתוצאה ממדיניות החוץ של ממשלת בריטניה, החליטו באוקטובר 1945 שלושת המחתרות לשתף פעולה באמצעות הקמת תנועת גג שתקרא ''תנועת המרי העברית``. 
	
	אופן הפעילות של תנועת המרי העברית, היה פגיעה צבאית ביעדים בריטיים, תשתיות בריטיות, וכו'. פרט למקרה יחיד, הפעילות הצבאית נעשתה ע''י מחתרת אחת בלבד (שכן המחתרות עדיין היו נפרדות), תוך ידיעה ותיאום עם שאר המחתרות. המקרה היחיד בו פעלו יחדיו היה ''ליל הרכבות`` – בלילה אחד פוצצו אנשי הפלמ''ח, זרוע צבאית מאומנת של ההגנה, ב־153 נקודות שונות במסילות הרכבת שח ארץ ישראל, ופגעו ב־3 ספינות של משמר החופים. בד בבד אנשי האצ''ל והלח''י תקפו את תחנת הרכבת המרכזית בלוד. 
	
	דוגמה אחרת לפעילות של תנועת המרי הוא ''ליל הגשרים`` – בלילה זה, ב־16 ביוני 1946, פוצצו 11 גשרים המקשרים בין א''י למדינות השכנות. שתי הדוגמאות מבטאות את אופי הפעילות של תנועת המרי העברי – פעילות צבאית כנגד מטרות ממשל בריטיות בולטות. 
	
	\section{}
	
	\textbf{תשובה: }\\
	בסוף המאה ה־19 התחלה להתגבש תנועת ה\textit{ציונות}, תנועה השואפת להקים בית מדיני לעם היהודי בארץ ישראל. התנועה החלה את פעולתה המאורגנת ע''י בינימין זאב הרצל, שהקים את רוב מסודותיה המאורגנים. בשנת 1945 נגמרה מלחמה בשם מלחמת העולם השנייה, שבה נכלאו ונעקרו מבתיהם יהודים רבים, שהיו צריכים למצוא מקום חדש לחיות בו. פתרון אפשרי לבעיה זו, שכנותה ''בעיית היהודים``, היה מעברם למגורים בא''י, בהתאם לרצונות התנועה הציונית. ואכן, רבים מן היהודים שעברו זוועות בזמן המלחמה, רצו וניסו להגדר לא''י. כל אותה התקופה, ועוד לפני המלחמה, פעלו מוסדות התנועה הציונית באופן דיפלומטי לשחרור הארץ מידי השלטון שהיה בה, ובאמצעים להעלת יהודים מאירופה. 
	
	בא''י באותה התקופה שלטה בריטניה, ונהגה במדיניות שנקראה ''ספר הלבן השלישי``. המדיניות הגבילה באופן משמעותי את כמות היהודים שיכלו לעלות לא''י, ואת התפשטות היישוב היהודי. פעלו בארץ שלוש מחתרות ציוניות עיקריות – ההגנה, שהייתה חלק מהנציגות ה''רשמית`` של התנועה הציונית, והובילה קו פעולה שהתאפיין בתמיכה בתחום הדיפלומטיה והעלייה, ואכן האצ''ל והלח''י שהתנגדו למדיניות ''ההגנה`` ופעלו נגד הבריטים, אזרחיהם ומוסדותיהם באופן צבאי. 
	
	ב־1945 לאחר סיום המלחמה עלתה לשלטון בבריטניה מפלגת הלייבור, שטרם עלייתה דמה כאילו היא תביא לתמיכה בציונים בארץ. עם זאת, שר החוץ של המפלגה, באוין, נעשה נאום שנקרא ''נאום באוין`` ובו הבהיר שמדיניות ממשלת הלייבור תהיה להמשיך לדבוק בספר הלבן השלישי. ביישוב היהודי, התחושה הייתה שאין להסתפק בדיפלומטיה, בהעפלה ועלייה בלבד, אלא שיש לפעול במקביל באופן צבאי. על כן, ההגנה, האצ''ל והלח''י החליטו לפעול באופן מתואם, והקימו באוקטובר 1945 את ''תנועת המרי העברי`` – תנועת גג שתתאם את פעילויות שלוש המחתרות העוינות. 
	
	לסיום, בעקבות נאום באוין והמדיניות של ממשלת הלייבור, שלושת המחתרות החליטו לשתף פעולה לאחר מלחמת העולם השנייה. 
	
	\section{}
	בסוף המאה ה־19 התחלה להתגבש תנועת ה\textit{ציונות}, תנועה השואפת להקים בית מדיני לעם היהודי בארץ ישראל. התנועה החלה את פעולתה המאורגנת ע''י בינימין זאב הרצל, שהקים את רוב מסודותיה המאורגנים. בשנת 1945 נגמרה מלחמה בשם מלחמת העולם השנייה, שבה נכלאו ונעקרו מבתיהם יהודים רבים. במלחמה שנארכה החל מ־1939 לחמו בפרט הבריטיים בגרמניה הנאצית, שפעלה כדי להשמיד ולרצוח את כל היהודים. בארץ ישראל פעלו שלוש מחתרות בזמן מלחמת העולם השנייה ועד 1948 – ההגנה, ארגון שניסה לפעול בעיקר באמצעים דיפלומטיים על־מנת להקים מדינה יהודית, וכן באמצעים של התיישבות ועליה (הבאת עולים לארץ ישראל, והקמת יישוב נוספים). בזמן המלחמה פרש מההגנה ארגון בשם האצ''ל, וממנו התפרק הלח''י בשלב מאוחר יותר. שני הארגונים האלו תמכו בפעילות צבאית נגד בריטניה ואזרחיה עוד במהלך המלחמה, תוך ידיעה שבריטניה לוחמת בגרמניה הנאצית שמשמידה את יהדות אירופה. 
	
	עוד משנת 1939 בריטניה פעלה במדיניות שנקראה ''הספר הלבן השלישי``. הספר הלבן השלישי הגביל אופן משמעותי את כמות היהודים שיכלו להתיישב בארץ, ואת מיקום ואופי היישובים החדשים שרצו הציונים לבנות. ממשלת הלייבור שעלתה בבריטניה בשנת 1945 ניפצה תקוות לסיום מדיניות הספר הלבן, בנאום שנשא שר החוץ של הממשלה – באוין – בו הוא הבהיר שהיהודים הרבים שנעקרו מבתיהם באירופה, אינם בעיה של ממשלת בריטניה, וכן זו האחרונה תמשיך במדיניות הספר הלבן השלישי. מדיניות החוץ של בריטניה הביאה את ההגנה ואת היישוב היהודים לחשיבה שלא די בפעילויות דיפלומטיות, עלייה והתיישבות, וכן יש צורך בפעילויות צבאיות. 
	
	התוצאה מחשיבה זו, ההגנה, האצ''ל והלח''י הסכימו לפעול באופן מתואם תחת ארגון גג שנקרא ''תנועת המרי העברי``, שהוקם באוקטובר 1945. תנועת המרי העברי פעלה במגוון צורות צבאיות כנגד השלטון הבריטי, בינהן פיצוץ פסי רכבת, הטבעת אוניות משמר החופים, פיצוץ גשרים, ועוד. למרות הפעולות הצבאיות המוצלחות, ביולי 1946 לאחר שמונה חודשים, התנועה התפרקה. נבחין בשתי סיבות עיקריות לכך: 
	\begin{itemize}
		\item התגובה הבריטית לפעילות תנועת המרי, הייתה תקיפה ביותר. נדגיש במיוחד את אירועי ''השבת השחורה`` – בזמן האירועים הבריטים תפסו כמויות ענק של נשק ברשות תנועת המרי, אסרו רבים מחברי התנועה, ותפסו עולים רבים. מבחינת היישוב היהודי וההגנה (אך לא האצ''ל והלח''י) השבת השחורה הוכיחה שהבריטים בעלי כוח רב בהרבה מהם, ושפעילות צבאית כנגדם רק תביא לפגיעה ביישוב היהודי. חוסר ההסכמה לגבי אופן הפעולה בין המחתרות השונות תרם לפירוק תנועת המרי. 
		\item סיבה אחרת היא פיצוץ מלון המלך דוד. אנשי האצ''ל גרמו לפיצוץ במלון בו בחלקו הדרומי יישבו חלקים מרחבים ממפקדי הצבא הבריטי. הפיצוץ אושר באופן עקרוני בתנועת המרי, אך בתנאי שלא יבוצע בשעות העבודה כדי להבטיח שלא תהיה פגיעה בנפש. האצ''ל פוצצו את חלקו הדרומי של המלון בשעות הצהריים בשיא פעולתו, בפיצוץ חזק שגם הרג אי אילו מהולכי הרגל (חלקם יהודים) שהסתובבו ברחוב יחד עם בלתי מעורבים רבים ששהו באיזור החלק הדרומי וכן בחלקו הדרומי של המלון. סך הכל נהרגו 91 אנשים, והאירוע גרם לנזק תדמיתי חמור ליישוב היהודי – האירוע נתפס ברחבי העולם כפיגוע טרור, וההגנה שרצתה בהקמת מדינה יהודית תוך שימוש בפעילות מדינית, גינתה בחריפות את אירוע. מסיבות אלו, פעילות משותפת של ההגנה והאצ''ל לא התאפשרה לאחר הפיצוץ, וזו סיבה נוספת לפירוק תנועת המרי. 
	\end{itemize}
	
	\section{}
	
	נערוך השוואה בין המקורות בעמוד 58: 
	
	\begin{tabularx}{\linewidth}{|X|l|l|l|l|c|}
		\hline תוכן & מחבר & זמן & קהל יעד & סוג &המקור  \\
		\hline גינוי חריף של התקיפה, תוך השמת דגש על חפותם של הנרצחים 
		 & עיתונאי ''הארץ`` & למחרת התקיפה & קוראי ''הארץ`` & כתוב & ראשון \\
		\hline טענה כי נתנו אזהרות טלפוניות רבות, ושאופי המאבק של האצ''ל לא פוגע באנשים, אלא אם זה נוגע לתלייה שלהם. לטענתו האשמה היא בבריטים שלא שעו לאזהרות ופינו את המלון & מפקד האצ''ל & 1978 & תושבי הארץ & ראיון & שני \\
		\hline טענה כי לא נמסר שום אזהרה טלפונית טרם הפיצוץ, ושגם אם הייתה נמסרת התראה כזו, פרק הזמן הקצר שעמד לרשות השוהים במלון לא היה די כדי לפנותם בבטחה. & מזכיר הבריטים & 1978 & תושבי הארץ & ראיון &שלישי \\
		\hline
	\end{tabularx}
	מההשוואה בין המקורות לפי הקריטיונים הבאים, נבחין שבעת התקיפה היישוב היהודי היה עסוק בגינוי התקיפה וביטול הנזק התדמיתי שנגרם, בעוד בשנים לאחריה העיסוק בה היה נוגע לעיצוב ההיסטוריה, ובעיקר עוסק סביב השאלה האם האצ''ל ניסה לפגוע באנשים, או לא (כלומר, נתן התראה טלפונית). 
	\section{}
	החל מסוף המאה ה־19 החלה להתגבש התנועה הציונית, שמנסה להקים מדינה לאומית על שטח ארץ ישראל. ראשית הארגונים הציוניים החלו בפועלו של בנימין זאב הרצל. בשנת 1939 השולת על א''י, בריטניה, הוציאה ספר חוקים בשם ''הספר הלבן השלישי`` שהגביל את עליית היהודים לארץ ישראל. גופים ציוניים בכללם ''ההגנה`` ו''המוסד לעלייה ב'`` פעלו ע''מ להביא עולים לארץ ישראל וליישבה בציוניים. שתי מטרות חשובות ניתן למנות להעפלה: 
	\begin{itemize}
		\item האחת, היא יישוב הארץ. רבים מחברי התנועה הציונית ראו ביישוב הארץ אמצעי שיעזור בעתיד להבהרה כי ארץ ישראל כוללות אוכלוסיה יהודית משמעותית ובתוכניות חלוקה עתידיות של הארץ, הציוניים יזכו בשטח. 
		\item השנייה, היא הפעלת לחץ דיפלומטי. ספינות מעפילים שלא צלחו בהגעה אל החוף, ונתקלו בכוחות בריטים שנקטו בכוח ע''מ להחזיר אותם לארצם או לגרש אותם לקפריסין, נראו ע''י מדינות העולם כאירועים תוקפניים ואנטי־לאומיים, והפעילו לחץ על הבריטים את מניעת ההעפלה וכן אשרור מדינה יהודית. 
	\end{itemize}
	נציג שתי פעולות שנעשו במסגרת ההעפלה. 
	\begin{itemize}
		\item הראשונה, היא פרשת לה־ספציה. 1אביב 1946 נעצצרו שתי ספינות מעפילים יהודים בנמל באיטליה, במטרה למנוע מהן להגיע לארץ. בתגובה, ראשי המוסד לעליה ב' ארגנו שביתת רעב של הנמצאים על הספינות שסוקרה בכל רחבי העולם. שביתת הרעב והלחץ על בריטניה אפשר לבסוף את מעבר אלפי המתיישבים לארץ ישראל,והטה את דעת הקהל הבין לאומי לטובת התנועה הציונית. 
		\item השני, ספינת המעפילים אקסודוס. על הספינה היו 4500 מעפילים, וספינות הצי הבריטי תקפו אותה בים (מחוץ למים הטריטוריאלים של בריטניה, ובניגוד לחוק הבינ''ל), התפתח קרב בו נהרגו מעפילים ונפצעו עשרות, והביטים החליטו להחזיר את האוניה לאירופה, שלאחר התנגדות בת 24 ימים לבסוף הועלו על ספינות והורדו באירופה. אומנם הספינה לא זכתה ליישב את ישראל, אך בכל הנוגע למטרה השנייה – מאות עיתונים סקרו את הפרשה והשפיעו על ועדות של האו''ם שעסקו בחלוקת הארץ. יתרה מכך – חברי הועדה שעסקה בבניית תוכנית חלוקה, נכחו במקום בו הורדו בכוח נוסעי אקסודוס. 
	\end{itemize}
	
	\section{}
	נתבונן בשני המקורות בעמוד 49. הראשון מביניהם מקור כתוב שמתאר חדשות, ועליו נכתב ''[ההגנה] זה שם ספינת המעפילים שהצליחה לפרוץ את ההסדר הבריטיה בחופי ההפלגה באירופה. הספינה מתקרבת חופי המולדת. יציאת אירופה תש''ז מסיעה לישראל: [...] 4554 נפש מישראל! חמאש אניות משחית ואניית סיור אחת שמו מצור על ספינת המעפילים לאורך כל הדרך. הצי הבריטי גייס כוח זה במטרה ללכוד את הספיהנ לבל יגיעו אחינו לחוף מבטחים``. נתייחס למטרות ההעפלה שתוארו בשאלה הקודמת, ונבחין כיצד האירועים המתוארים במקור זה מסייעים למטרות ההעפלה: 
	\begin{itemize}
		\item האירוע הראשון, הוא הבקעת הספינה ''ההגנה`` את המחסומים הבריטים. בכך, הנמצאים על הסיפון יוכלו לעלות לארץ, ולעזור ליישוב הארץ. 
		\item האירוע השני, הוא אסירתם של כ־4.5 אלף אנשים, וגירושם. גירושם עוזר  לגיוס לדעת הקהל הבין לאומית לצד התנועה הציונית. 
		\item השלישי, הוא שחמש אוניות משחית ואונית סיור שמו מצור על ספינת מעפילים. האירוע הזה גם עוזר לגייס את דעת הקהל הבינ''ל. 
	\end{itemize}
	
	נתבונן עכשיו במקור השני – מקור ויזואלי, צילום של עיתונאי המתאר את גירוש אקסודוס. הצילום מראה ספינה הרוסה כליל, ממנה יוצאים אנשים בעצב רב, חיילים בריטים מאגפים אותם. כמתואר בשאלה הקודמת, תמונות מעין אלו שעיתונאים הפיצו עזרו לגייס את דעת הקהל הבינ''ל לצד התנועה הציונית. 
	
	\section{}
	החל מסוף המאה ה־19 החלה להתגבש התנועה הציונית, שמנסה להקים מדינה לאומית על שטח ארץ ישראל. ראשית הארגונים הציוניים החלו בפועלו של בנימין זאב הרצל. בשנת 1939 השולת על א''י, בריטניה, הוציאה ספר חוקים בשם ''הספר הלבן השלישי`` שהגביל את התרחבות היישוב היהודי. גופים ציוניים בכללם ''ההגנה`` ו''קק''ל`` (כולם חלק מההסתדרות הציונית) פעלו על מנת ליישב את הארץ בציונים. שתי מטרות חשובות ניתן להתיישבות בארץ: 
	\begin{itemize}
		\item חלק ניכר מהיישובים הוקמו במגמה התיישבותית שנקראה ''חומה ומגדל`` – יישוב שנבנו באוזרים אסטרטגים ע''מ לעזור לבלום פלישה של צבאות ערב, שהיה ברור לכל שתתרחש במידה ותקום מדינה יהודית. ההתיישבות עזרה כך לחזק את בטחון המדינה הצפויה. 
		\item היישובים נבחרו במקום אסטרטגיים, כך שהיישוב היהודי יתפרש על שטח נרחב מא''י ככל האפשר – כך, קיוו הציונים, שכאשר העולם יכריע את גבולותיה של מדינה עתידית, חלק ניכר מהשטח יהיה שייך למדינה היהודית. 
	\end{itemize}
	
	נתבונן בשתי פעולות בתחום הזה, ונסביר כיצד הן קידמו את המטרות לעיל: 
	\begin{itemize}
		\item הראשון, הוא מבצע י''א הנקודות. בלילה אחד, במוצאי יום הכיפורים 1946, 11 יישובים בנגב הוקמו ללא ידיעת הבריטים. במבצע היו מעורכים 300 משאיות של חברת ''סולל בונה``, כ־1100 איש, אנשי אבטחה, ותמיכה נרחבת של ארגונים כמו ''ההגנה`` והחברה ''מקורות`` (ספקית המים של היישוב היהודי). השלטונות הבריטים אשרו בדיעבד את הקמת היישובים. גוש היישובים יצר אחזיה יהודית בנגב, והוביל, לאחר מכן, להכנסת הנגב למדינה היהודית בתוכנית החלוקה של שנת 1947. 
		\item פרשה אחרת שנדגים הייתה ''פרשת ביריה``. ב־28 בפברואר הצבא הבריטי מצא נשק בסמוך למאחז היהודי ביריה, ולאחר שמצא שם נשק, גירש את תושבי ביריה והחליפם ביחידה צבאית בריטית. מחשש לפעולות נוספות דומות, למעלה מ־2000 בני נוער עשו עצמם כעולים לתל חי, ובפועל הקימו יישוב חדש במקום ביריה. אחר פינוי המקום, מאות אנשים זרמו אליו ולא אפשרו את הריסתו או כיבושו. לבסוף ביריה נשאר יישוב עברי. בכך, הפעולה הצליחה למנוע את נפילתם של יישובים נוספים לשלטון הבריטי, ובכך לקדם את כל מטרות ההתיישבות. 
	\end{itemize}
	
	\section{}
	במהלך המאה ה־17 ועד המאה ה־20 בריטניה עסקה בתהליך שנקרא קולוניזציה – הקמת קולוניות ברחבי העולם, במטרה לנצל את משאביהם (הטבעיים והאנושיים) ולחזק את שליטה הצבאית בעולם. החל מסוף המאה ה־19 החלה להתגבש התנועה הציונית, שמנסה להקים מדינה לאומית על שטח ארץ ישראל. ראשית הארגונים הציוניים החלו בפועלו של בנימין זאב הרצל. בשנים 1939-1945 התרחשה ''מלחמת העולם השנייה``, בה בריטניה לחמה כנגד מדינות כמו גרמניה הנאצית ואיטליה הפאסישטית. למרות הנצחון במלחמה, המלחמה רוששה את בריטניה, וזו הפכה להיות תלויה כלכלית בארצות הברית. בשנת 1945 הייתה ישראל קולוניה של בריטניה. ישראל הייתה יעד חשוב – האיזור הגיאוגרפי בעל השפעה מכרעת (גובל בים האדום ובים התיכון, סמוך לתעלת סואץ, והיבשה היחידה שמפרידה את כל אסיה מכל אפריקה). ישראל עזרה לישראל להעביר סחורות מהודו, קולוניה אחרת בשליטתה, ולהעמיק את האיחזותה במזרח התיכון. 
	
	המצב הכלכלי בבריטניה בסוף המלחמה היה בכי רע, ונסיונם לשלוט בקולניות בכל העולם עלה להם מחיר כלכלי כבד – היה חייל בריטי כנגד כל פחות מ־20 אנשים בא''י, והיישוב הציוני נאבק באופן מתמשך בשלטון הבריטי. מלחמות חוזרות ונשנות עלו בין היהודים לערבים שחיו באיזור. מעצמות וקהל העולם ראו ביישוב הציוני כצד הצודק והמוסרי, וארה''ב הפעילה לחצים רבים על בריטניה לאפשר עלייה של כמאה אלף יהודים. 
	
	בסיום מלחמת העולם השנייה הוקם האו''ם, גוף בין־לאומי שמטרותו לשמור על השלום ולתאם בין המדינות השונות, במטרה למנוע הישנות של מלחמה כמו מלחמת העולם השנייה בעתיד. 
	
	נוכל למנות מספר סיבות מדוע בשנת 1947 בריטניה העבירה את ההחלטה בדבר ישראל לאו''ם. 
	
	\begin{itemize}
		\item בריטניה שהייתה תלויה כלכלית בארה''ב הייתה בלחץ רב לקבל את הצעתם לאפשר עלייה של כמאה אלף יהודים. מצד שני, היה ברור שהמלך ייצור מלחמות רבות בארץ בין היהודים לערבים, ויגרום להרעת מצב בריטניה ביחס לכל מדינות הערב במזרח התיכון. 
		\item היו קשיים כלכליים רבים בבריטניה, ולא השתלם לה כלכלית להחזיק את השטח. 
		\item בריטניה התכוננה לצאת מהודו, כלומר נעלמה החשיבות של ישראל כמקום להעברת סחורות דרכו. 
		\item דעות הקהל העולמי נטו נגד בריטניה. 
	\end{itemize}
	מסיבות אלו, בריטניה החליטה להעביר את ההחלטה לאו''ם, מה שהוביל ליציאתה מישראל. בכך, בריטניה לא פתחה חזית נגד הערבים, והוקל עליה כלכלית שכן נעלם הצורך להחזיק בא''י. 
	
	\ndoc
\end{document}